

The WIMP signal considered in this analysis is expected to produce low-energy, single-scatter NR signals uniformly distributed in the TPC, with no additional signals in the TPC, skin, or OD. 
The following strategy is used to obtain a clean sample of such events: exclude time periods of elevated TPC activity or electronics interference; remove multiscatter interactions in the TPC; remove events outside an energy region of interest (ROI); remove events due to accidental coincidence of S1 and S2 pulses; remove events with coincident signals in the TPC and skin or OD; remove events near the TPC active volume boundaries.
Methods of bias mitigation that involve obscuring the data, such as blinding the signal region or adding fake events (``salting''), were avoided to allow control over larger sources of systematic errors that may be presented by a new detector. To mitigate bias in this result, all analysis cuts were developed and optimized on sideband selections and calibration data. 




The search dataset totals 89~live days after removing periods for detector maintenance and calibration activity, as well as a \SI{3}{\percent} loss due to DAQ dead time and a \SI{7}{\percent} loss to periods excised due to anomalous trigger rates. Because dual-phase xenon TPCs experience elevated rates of activity after large S2 pulses~\cite{akerib2017results, collaboration2018dark,sorensen2017electron, LUXelectronbkgs}, a time hold-off is imposed to remove data taken after large S2s and after cosmic-ray muons traversing the TPC. These omissions result in a final search live time of \SI{60(1)}{d} where a WIMP interaction could be reconstructed. 
In future searches, the hold-off can be relaxed by optimization with respect to analysis cuts and detector operating conditions.

The ROI is defined as \sonec~in the range $3-80$\,phd, uncorrected S2 greater than \SI{600}{phd} ($>$10 extracted electrons), and \stwoc~less than $10^5$\,phd, ensuring that signal efficiencies are well understood and background ER sources are well calibrated by the tritium data.
Events classified as multiple scatters in the TPC are removed, as are events with poor reconstruction due to noise, spurious pulses, or other data anomalies.

A suite of analysis cuts targets accidental coincidence events, henceforth called ``accidentals,'', where an isolated S1 and an isolated S2 are accidentally paired to mimic a physical single-scatter event. Isolated S1s can be generated from sources such as particle interactions in charge-insensitive regions of the TPC, Cherenkov and fluorescent light in detector materials, or dark-noise pileup. Isolated S2s can be generated from sources such as radioactivity or electron emission from the cathode or gate electrodes, particle interactions in the gas phase or in the liquid above the gate electrode, or drifting electrons trapped on impurities and released with $\mathcal{O}$(\SI{100}{ms}) time delay \cite{LUXelectronbkgs}.
Analysis cuts to remove accidentals target individual sources of isolated S1s and S2s using the expected behavior of the S1 and S2 pulses with respect to quantities such as drift time, top-bottom asymmetry of light, pulse width, timing of PMT hits within the pulse, and hit pattern of the photons in the PMT arrays. The cuts remove $>$\SI{99.5}{\percent} of accidentals, measured using single-scatter-like events with unphysical ($>$\SI{951}{\micro\second}) drift time and events generated by random matching of isolated S1 and S2 populations. 

\begin{figure}[htbp]
  \centering
  \includegraphics[ width=\linewidth]{figures/FlatNR_efficiencies_NoSkew_ERTune_S2RecoEff.pdf}
  \caption{Signal efficiency as a function of NR energy for the trigger (blue), the threefold coincidence and $>$\SI{3}{phd} threshold on S1c (orange), single-scatter (SS) reconstruction and analysis cuts (green), and the search ROI in S1 and S2 (black). The inset shows the low-energy behavior, with the dotted line at \SI{5.5}{keV_{\text{nr}}} marking  \SI{50}{\percent} efficiency. 
  The error band (gray) is assessed using AmLi and \tritium\ data as discussed in the text.
  }
  \label{fig:NRacc}
\end{figure}


Data-driven signal efficiencies for the trigger, reconstruction, and analysis cuts are shown in Fig.~\ref{fig:NRacc}. The DAQ trigger efficiency is determined from DD data by comparing the external trigger of the generator against the TPC S2 trigger logic, and is confirmed using randomly triggered events collected throughout the search.
The reconstruction efficiency for low-energy NR events is evaluated by comparing the reconstruction results against a large set of events manually identified as single scatter in DD data. An additional reconstruction inefficiency due to S2 splitting for long drift times for low numbers of extracted electrons is accounted for with simulation.
Analysis cut efficiency is not determined directly from neutron calibration data as they do not cover the spatial extent of the TPC and are contaminated by a high rate of single photons and electrons. 
Instead, the efficiency throughout the full analysis volume is evaluated using \tritium\ data for analysis cuts targeting S1 pulses and the combination of \tritium\ and AmLi data for those targeting S2 pulses. Composite NR-like waveforms are generated using \tritium\ single scatters with their S2 pulses replaced by smaller pulses from other \tritium\ or AmLi events (an ``AmLi-\tritium'' dataset).  The uncertainty on the NR signal efficiency is the larger of the $\pm1\sigma$ statistical fluctuation of the AmLi-\tritium\ dataset and the difference between the AmLi-\tritium\ dataset and a pure AmLi dataset. The uncertainty is \SI{3}{\percent} for nuclear recoil energies $>$\SI{3.5}{keV_{\text{nr}}}, increasing to \SI{15}{\percent} at \SI{1}{keV_{\text{nr}}}.

Events with coincident activity in the TPC and skin or OD are removed to reduce backgrounds producing $\gamma$ rays and neutrons. To mitigate backgrounds associated with $\gamma$ rays, events with a prompt signal in the OD (skin) within $\pm$\SI{0.3}{\micro\s} ($\pm$\SI{0.5}{\micro\s}) of the TPC S1 pulse are removed. Neutrons can thermalize in detector materials and those that capture on  hydrogen or gadolinium in the OD can be tagged by an OD pulse of greater than $\sim$\SI{200}{keV} within \SI{1200}{\micro\s} after the TPC S1. A selection on large skin pulses in the same time window additionally tags $\gamma$ rays returning to the xenon from an OD capture process. AmLi calibration sources placed at the nine locations close to the TPC are used to determine a position-averaged neutron tagging efficiency of \SI{89(3)}{\percent} for TPC single-scatters in the nuclear recoil band. Background data is used to determine a false veto rate of $5\%$ due to accidental activity in the OD during the coincidence window. Background neutrons may have a higher tagging efficiency due to their harder energy spectrum and coincident $\gamma$-ray emission.

\begin{figure}[tbp]
 \centering
 \includegraphics[ width=\linewidth]{figures/SR1WS_rdrift_OD_SKIN_20220705.pdf}
 \caption{Data in reconstructed $r^2$ and $z$ after all analysis cuts. Black (gray) points show the data inside (outside) the FV. Red crosses and blue circles show events vetoed by a prompt LXe skin or OD signal, respectively. The solid line shows the FV definition, and the dashed line shows the extent of the active TPC. Field nonuniformities cause the reconstructed $r$ position of the active volume boundary to vary as a function of z. Events with drift time of approximately \SI{50}{\micro\s} are from recoils in the gas that produce S1 and S2 pulses with a fixed time separation.}
 \label{fig:r2z}
\end{figure}

Finally, events outside a central fiducial volume (FV) are removed to reject external and other backgrounds that concentrate near the TPC boundaries, as shown in Fig.~\ref{fig:r2z}. Events at high radius have reduced position reconstruction resolution due to reduced S2 light collection efficiency and charge-loss effects within a few millimeters of the polytetrafluoroethylene wall.
The radial extent of the FV and the S2 threshold are chosen simultaneously using data outside the \sonec\ ROI to eliminate events leaking into the FV due to poor position reconstruction resolution. Radially, the FV terminates at \SI{4.0}{cm} in reconstructed position from the TPC wall, with small additional volumes removed in the top (\SI{5.2}{cm} for drift time $<$\SI{200}{\micro\s}) and bottom (\SI{5.0}{cm} for drift time $>$\SI{800}{\micro\s}) corners to account for increased rates of background in those locations. Events within \SI{6.0}{cm} of the $(x, y)$ positions of two ladders of TPC field-cage resistors embedded in the TPC wall are also removed.  Vertically, events with drift times $<$\SI{86}{\micro\s} and $>$\SI{936.5}{\micro\s} are rejected, corresponding to \SI{12.8}{cm} and \SI{2.2}{cm} from the gate and cathode electrodes, respectively.
The number of remaining events from the wall entering the FV is estimated to be $<0.01$.
The xenon mass in the FV is estimated to be \SI{5.5(2)}{\tonne} using tritium data and confirmed by geometric calculation. 

\begin{figure}[tbp]
  \centering
  \includegraphics[ width=\linewidth]{figures/s1logs2_SR1_WS_30GeV_revised.pdf}
  \caption{WIMP-search data (black points) after all cuts in $\log_{10}$\stwoc-\sonec~space. Contours enclose $1\sigma$ and $2\sigma$ of the following models: the best-fit background model (shaded gray regions), the \ArThreeSeven\ component (orange ellipses), a \SI{30}{GeV/c\squared} WIMP (purple dashed lines), and $^8$B solar neutrinos (shaded green regions).
  The red solid line indicates the NR median, and the red dotted lines indicate the \SI{10}{\percent} and \SI{90}{\percent} quantiles.
  Model contours incorporate all efficiencies used in the analysis. 
  Thin gray lines indicate contours of constant energy.
  }
  \label{fig:s1s2wimp}
\end{figure}

Figure~\ref{fig:s1s2wimp} shows the distribution in $\log_{10}$\stwoc{\nobreakdash-}\sonec\ of the 335 events~\footnote{See Supplemental Material %at \url{\supplementalurl} 
for a detailed table of events by selection.}  
passing all selections, along with contours representing a \SI{30}{GeV/c\squared} WIMP, a flat NR distribution, and the background model.
The signal model assumes spin-independent scattering from WIMPs with an isotropic Maxwell–Boltzmann velocity distribution, parameterized as in Ref.~\cite{baxter2021recommended}, with inputs from Refs.~\cite{LEWIN199687,smith2007rave,mccabe2014earth, schonrich2010local,blandhawthorn2016galaxy, gravity2021improved}.
The WIMP model has an approximately exponentially decreasing energy spectrum with shape that depends on the mass of the WIMP~\cite{LEWIN199687}.


